% Zdroj: PDF priklady-podzim-2024.pdf, fyzická str. 1-2. Značka: společné okruhy. Zařazeno: I1.
% Obrázky: 3 stavové diagramy oříznuty z PDF (NFA A + DFA B ze str. 1, doplňkový DFA C ze str. 2).
\section{Doplněk regulárního výrazu}
\szzotazka{podzim 2024}{1}{Doplněk regulárního výrazu (NFA/DFA, podmnožinová konstrukce, doplněk)}

\noindent\textbf{Zadání.}\par
Mějme následující regulární výraz nad abecedou $\Sigma = \{a, b\}$:
\[ R = ((a+b)(a+b))^* ab. \]
Označme jako $L$ regulární jazyk popsaný výrazem $R$.
\begin{enumerate}
\item Sestrojte (co nejmenší) \textit{nedeterministický} konečný automat $A$
  rozpoznávající jazyk $L$.
\item Pomocí podmnožinové konstrukce převeďte automat $A$ na \textit{deterministický}
  konečný automat $B$.
\item Z~automatu $B$ sestrojte \textit{deterministický} konečný automat $C$
  rozpoznávající \textit{doplněk} jazyka $L$.
\end{enumerate}
Sestrojené automaty $A$, $B$ a~$C$ znázorněte pomocí stavových diagramů.

\medskip
\noindent\textbf{Náčrt řešení.}\par
\begin{enumerate}
\item Nejmenší NFA má 4 stavy: $A = (\{q_0, q_1, q_2, q_F\}, \{a, b\}, \delta_A, q_0,
  \{q_F\})$, kde $\delta_A$ popíšeme stavovým diagramem:
  \begin{center}\includegraphics[width=0.75\linewidth]{doplnek-nfa-a.png}\end{center}

\emergencystretch=1.5em
\item Výsledkem podmnožinové konstrukce je následující DFA: $B = (\{\{q_0\},\allowbreak \{q_1\},\allowbreak
  \{q_1, q_2\},\allowbreak \{q_0, q_F\}\},\allowbreak \{a, b\},\allowbreak \delta_B, q_0, \{\{q_0, q_F\}\})$, kde
  $\delta_B$ popíšeme stavovým diagramem:
  \begin{center}\includegraphics[width=0.8\linewidth]{doplnek-dfa-b.png}\end{center}

\item Stačí zaměnit přijímající a~nepřijímající stavy:
  \begin{center}\includegraphics[width=0.8\linewidth]{doplnek-dfa-c.png}\end{center}
  (Všimněte si, že přechodová funkce automatu $B$ je totální, proto není třeba
  přidávat FAIL stav.)
\end{enumerate}

\medskip
\noindent\textit{Doplňující vysvětlení (nad rámec oficiálního náčrtu):} $L$ tvoří
slova sudé délky končící na $ab$. Stav $q_0$ znamená „přečten sudý počet znaků``,
$q_1$ „lichý počet``, z~$q_0$ lze na $a$ nedeterministicky zahájit závěrečné $ab$
(přes $q_2$ do $q_F$). Méně než 4 stavy NFA mít nemůže: pro dvojice slov
$(x_i, y_i) \in \{(\varepsilon, ab), (a, b), (ab, \varepsilon), (b, aab)\}$ je vždy
$x_i y_i \in L$, ale pro $i \neq j$ do $L$ nepatří $x_i y_j$ nebo $x_j y_i$. Stavy, ve
kterých jsou přijímající výpočty slov $x_i y_i$ po přečtení $x_i$, proto musí být
navzájem různé.
